\chapter{Manual de Usuario e Instalación}
\label{ape:manual}

Este anexo proporciona las instrucciones necesarias para la puesta en marcha y operación del sistema en un entorno local, incluyendo los requisitos de hardware y software, el procedimiento de instalación paso a paso, la guía de operación con ejemplos reales y una sección de resolución de problemas comunes.

\section{Requisitos del Sistema}

\subsection{Hardware recomendado}
Para una ejecución fluida de los modelos locales, se recomienda:
\begin{itemize}
    \item \textbf{GPU:} NVIDIA con arquitectura CUDA (mínimo 8 GB VRAM, ej. RTX 3060/4050 o superior). Sin GPU, el sistema opera por CPU con tiempos de ingesta significativamente mayores.
    \item \textbf{RAM:} 16 GB o superior. El sistema mantiene los grafos de todos los documentos cargados en memoria durante la sesión (aproximadamente 15 MB por cada 100 páginas procesadas).
    \item \textbf{Almacenamiento:} Mínimo 10 GB de espacio libre para alojar los modelos de lenguaje descargados por \texttt{Ollama} y los modelos de \textit{embeddings} de \texttt{sentence-transformers}.
\end{itemize}

\subsection{Software necesario}
\begin{itemize}
    \item \textbf{Python 3.10} o superior.
    \item \textbf{Ollama:} Motor de inferencia local. Disponible en \url{https://ollama.com}. Debe estar en ejecución como servicio en segundo plano antes de iniciar el sistema.
    \item \textbf{NVIDIA CUDA Toolkit:} Versión 12.x para aceleración por hardware. Descargable desde \url{https://developer.nvidia.com/cuda-downloads}.
    \item \textbf{Git:} Para clonar el repositorio del proyecto.
\end{itemize}

\section{Instalación y Configuración}

\subsection{Paso 1: Obtención del código fuente}
Clone el repositorio del proyecto en su directorio de trabajo:
\begin{verbatim}
git clone <URL_del_repositorio>
cd <nombre_del_directorio>
\end{verbatim}
Sustituya \texttt{<URL\_del\_repositorio>} por la URL real del repositorio.

\subsection{Paso 2: Preparación del entorno virtual}
Se recomienda el uso de un entorno virtual de Python para aislar las dependencias del proyecto:
\begin{verbatim}
python -m venv venv

# En Linux/macOS:
source venv/bin/activate

# En Windows:
venv\Scripts\activate
\end{verbatim}
Una vez activado el entorno, instale todas las dependencias del proyecto:
\begin{verbatim}
pip install -r requirements.txt
\end{verbatim}

\subsection{Paso 3: Descarga de modelos en Ollama}
Con \texttt{Ollama} instalado y en ejecución, descargue el modelo de lenguaje principal:
\begin{verbatim}
ollama pull llama3.1:8b
\end{verbatim}
Este proceso requiere una conexión a internet y descarga aproximadamente 4,7 GB. Los modelos de \textit{embeddings} (\texttt{sentence-transformers/all-MiniLM-L6-v2} y el \textit{cross-encoder}) se descargan automáticamente desde HuggingFace en la primera ejecución del sistema.

\subsection{Paso 4: Preparación de los documentos}
Coloque los archivos PDF que desea analizar en el directorio \texttt{Code/MVP/Pliegos/} del proyecto (o modifique la variable de ruta al inicio de \texttt{main.py} para apuntar a su directorio de documentos):
\begin{verbatim}
mkdir Code/MVP/Pliegos
cp /ruta/a/sus/pliegos/*.pdf Code/MVP/Pliegos/
\end{verbatim}

\section{Guía de Operación}

\subsection{Inicio del sistema}
Asegúrese de que \texttt{Ollama} esté en ejecución en segundo plano y, con el entorno virtual activado, ejecute el script principal:
\begin{verbatim}
python main.py
\end{verbatim}
Al iniciarse, el sistema procesará automáticamente todos los PDFs del directorio configurado. Verá en la consola el progreso de extracción, generación de \textit{embeddings} y construcción del grafo para cada documento. Una vez finalizada la fase de ingesta, aparecerá el \textit{prompt} interactivo de chat.

\subsection{Comandos de control}
Durante la sesión de chat, el usuario puede utilizar los siguientes comandos especiales para gestionar el comportamiento del sistema:

\begin{table}[ht]
\centering
\begin{tabularx}{\textwidth}{|l|X|}
\hline
\textbf{Comando} & \textbf{Efecto} \\ \hline
\texttt{modo 1} & Cambia a respuestas individuales por documento. Cada PDF recibe y responde la consulta de forma independiente. Modo por defecto. \\ \hline
\texttt{modo 2} & Activa el modo consolidado multidocumento. El sistema extrae subgrafos de todos los documentos y genera una única respuesta integrada. \\ \hline
\texttt{visualizar si} & Activa la generación de archivos \texttt{.png} con la representación visual del subgrafo de consulta tras cada respuesta. \\ \hline
\texttt{visualizar no} & Desactiva la generación de visualizaciones (comportamiento por defecto). \\ \hline
\texttt{salir} & Finaliza la ejecución del programa de forma limpia. \\ \hline
\end{tabularx}
\caption{Comandos de control disponibles en la interfaz interactiva.}
\label{tab:comandos}
\end{table}

\subsection{Ejemplos de consultas}

A continuación se presentan ejemplos representativos de consultas y el tipo de respuesta que genera el sistema:

\textbf{Consulta específica (modo PARTIAL):}

\medskip\noindent\textit{¿Qué medidas de seguridad física exige el pliego para el rack de comunicaciones?}

El agente de estrategia selecciona \texttt{PARTIAL}, el optimizador reformula la consulta a \textit{``seguridad física rack comunicaciones''}, y el sistema recupera el nodo correspondiente junto con su sección padre. La respuesta detalla los requisitos de anclaje, cierre y acceso restringido del \textit{rack}.

\textbf{Consulta de síntesis global (modo FULL):}

\medskip\noindent\textit{Resume los principales requisitos técnicos del sistema de videovigilancia.}

El agente detecta la necesidad de síntesis y selecciona \texttt{FULL}, enviando al LLM el contenido completo del documento. La respuesta consolida los requisitos de resolución de cámara, almacenamiento y conectividad en una lista estructurada.

\textbf{Consulta multidocumento (modo 2):}

\medskip\noindent\textit{Compara los plazos de entrega del pliego técnico y el administrativo.}

Con el modo 2 activo, el sistema consulta ambos grafos de forma simultánea y genera una respuesta que cita explícitamente cada documento, señalando posibles inconsistencias entre ellos.

\section{Resolución de Problemas Frecuentes}

\begin{description}
    \item[El sistema no detecta secciones del índice correctamente.] Este problema suele darse en documentos cuyo índice no contiene hipervínculos internos o cuya numeración de secciones es irregular. El sistema intentará la extracción mediante expresiones regulares como mecanismo de \textit{fallback}. Si los resultados siguen siendo incorrectos, verifique que el PDF no esté protegido contra la extracción de texto o que no sea un documento escaneado sin capa OCR.

    \item[Ollama no responde o lanza un error de conexión.] Compruebe que el servicio de \texttt{Ollama} esté en ejecución ejecutando \texttt{ollama list} desde la terminal. Si el servicio no aparece, inícielo con \texttt{ollama serve}. Verifique también que el puerto por defecto (11434) no esté bloqueado por el cortafuegos del sistema.

    \item[Error de memoria insuficiente (CUDA out of memory).] El modelo \texttt{llama3.1:8b} cuantizado a 4-bit requiere aproximadamente 5,2 GB de VRAM. Si su GPU dispone de menos memoria, puede intentar usar el modelo en su variante de menor tamaño (\texttt{llama3.2:3b}) ejecutando \texttt{ollama pull llama3.2:3b} y configurando el nombre del modelo en el código fuente.

    \item[Los tiempos de ingesta son muy elevados.] Si el sistema opera en modo CPU (sin GPU NVIDIA o sin CUDA instalado), los tiempos de generación de \textit{embeddings} pueden ser entre $3\times$ y $5\times$ superiores. Asegúrese de que \texttt{CUDA Toolkit 12.x} esté correctamente instalado y de que \texttt{torch} detecte la GPU ejecutando:
\begin{verbatim}
python -c "import torch; print(torch.cuda.is_available())"
\end{verbatim}

    \item[El sistema genera un chunk que mezcla contenido de secciones distintas.] Esto ocurre ocasionalmente cuando los títulos de sección tienen una similitud textual muy alta y el \textit{fuzzy matching} no puede distinguirlos con precisión. En ese caso, puede ajustar el umbral de similitud difusa en el archivo de configuración de \texttt{chunk\_extractor.py}.
\end{description}
